% !TeX program = pdflatex
% =============================================================================
% Experimente 1: Normalkraft
% UZH Praktikum I -- Sara Romer Urban, Kantonsschule im Lee, HS2026
% Klasse: 2e (09.09.2026)
%
% Ein kurzer Demoversuch als Ergaenzung zu normalkraft-lektionsplan.tex (nicht
% Ersatz fuer Arbeitsblatt/Aufgabenblatt): Brett-Biege-Versuch -- ein Brett auf
% zwei Stuetzen biegt sich sichtbar durch, sobald eine Masse mit ihrer
% Gewichtskraft $F_g$ in die Mitte gelegt wird (nicht eine Person, die auf das
% Brett steigt -- wird im Unterricht live mit einem echten Brett und einer
% echten Masse gezeigt). Zeigt, dass auch eine scheinbar starre Oberflaeche
% elastisch nachgibt (Verallgemeinerung der Tisch-Idee aus dem Arbeitsblatt auf
% ein Objekt, das sich tatsaechlich sichtbar verformen laesst).
% =============================================================================
\documentclass[11pt,a4paper]{article}

\usepackage[T1]{fontenc}
\usepackage[utf8]{inputenc}
\usepackage[ngerman]{babel}
\usepackage{lmodern}
\usepackage[a4paper,margin=2.2cm]{geometry}
\usepackage{amsmath}
\usepackage{siunitx}
\sisetup{per-mode=fraction,fraction-command=\frac}
\usepackage{xcolor}
\usepackage{tikz}
\usepackage{graphicx}
\usetikzlibrary{arrows.meta,calc}

\definecolor{titleblue}{HTML}{183A6B}
\definecolor{accentorange}{HTML}{D9720C}
\definecolor{accentblue}{HTML}{2E6F9E}
\definecolor{darkgray}{HTML}{4A4A4A}
\definecolor{warnred}{HTML}{9E2E2E}

\setlength{\parindent}{0pt}
\setlength{\parskip}{0.6em}

\usepackage{titlesec}
\titleformat{\section}{\Large\bfseries\color{titleblue}}{\thesection}{0.6em}{}
\titlespacing*{\section}{0pt}{1.2em}{0.5em}

\begin{document}
\sloppy

\begin{center}
{\Huge\bfseries\color{titleblue}Experimente 1: Normalkraft}\\[0.4em]
{\large Kurze Demoversuche zur Doppellektion}
\end{center}
\vspace{0.5em}

\noindent\textbf{Klasse:} 2e (09.09.2026) \hfill \textbf{Lehrperson:} Loris Keller

\vspace{0.3em}
\noindent\textcolor{darkgray}{Dieses Dokument sammelt kurze Demoversuche zur
Doppellektion Normalkraft, mit Skizze und kurzer Erklärung pro Versuch, als
Ergänzung zum Unterricht, nicht als Ersatz für Arbeitsblatt oder
Aufgabenblatt.}

% =============================================================================
\section*{Experiment: Biegt sich auch ein starres Brett?}

\textbf{Aufbau:} Ein Holzbrett wird auf zwei Stützen (z.\,B. Bücherstapel
oder Hocker) gelegt, sodass es in der Mitte frei liegt. Anschliessend
wird eine Masse mit der Gewichtskraft $\vec F_g$ in die Mitte des
Bretts gelegt.

\begin{center}
\begin{tikzpicture}[>=Latex,scale=1.1]
  \draw[very thick,darkgray] (0,0) rectangle (1.2,1.6);
  \draw[very thick,darkgray] (6.8,0) rectangle (8,1.6);
  \draw[thick,darkgray] (0,1.6) -- (8,1.6);
  \draw[fill=lightgray!60] (3.6,1.6) rectangle (4.4,2.2);
  \node at (4,1.9) {\scriptsize $m$};
  \node[below] at (4,-0.3) {\small Vorher: Brett gerade};
\end{tikzpicture}
\hspace{0.6cm}
\begin{tikzpicture}[>=Latex,scale=1.1]
  \draw[very thick,darkgray] (0,0) rectangle (1.2,1.6);
  \draw[very thick,darkgray] (6.8,0) rectangle (8,1.6);
  \draw[thick,darkgray] (0,1.6) .. controls (3,1.6) and (3,0.7) .. (4,0.7)
    .. controls (5,0.7) and (5,1.6) .. (8,1.6);
  \draw[fill=lightgray!60] (3.6,0.7) rectangle (4.4,1.3);
  \node at (4,1.0) {\scriptsize $m$};
  \fill (4,1.0) circle (0.5pt);
  \draw[->,ultra thick,accentorange] (4,1.0) -- (4,-0.1) node[midway,right]{$\vec F_g$};
  \node[below] at (4,-0.4) {\small Nachher: Brett durchgebogen};
\end{tikzpicture}
\end{center}

\textbf{Erklärung:} Auch ein Brett, das auf den ersten Blick völlig
starr wirkt, verformt sich unter Belastung sichtbar, genau wie die
weiche Unterlage aus dem Arbeitsblatt, nur deutlich weniger stark. Die
Verformung ist der sichtbare Beweis dafür, dass das Brett eine
Normalkraft $F_N$ nach oben aufbaut, die $\vec F_g$ das Gleichgewicht
hält. Ein normaler Tisch verformt sich nach demselben Prinzip, nur so
wenig, dass es mit blossem Auge nicht mehr sichtbar ist.

\end{document}
